% Part: normal-modal-logic
% Chapter: tableaux
% Section: rules-for-K

\documentclass[../../../include/open-logic-section]{subfiles}

\begin{document}

\olfileid{nml}{tab}{rul}

\olsection{Aturan kanggo \Ax{K}}

Aturan kanggo panggandheng proposisional lumrah padha karo aturan
kanggo !!{tableau}s proposisional mawa tandha lumrah, mung ditambahi
prefiks. Ing saben kasus, aturan kang ditrapake marang !!{formula}
mawa tandha $\sFmla{S}{!A}[\sigma]$ ngasilake !!{formula}s anyar kang
uga nduweni prefiks~$\sigma$. Iki cetha miturut intuisi: upamane, yen $!A
\land !B$ bener ing jagad kang dijenengi~$\sigma$, mula $!A$ lan $!B$
bener ing~$\sigma$ (dene aturan iki ora mratelakake kabenerane ing jagad
liyane). Aturan
proposisional diklumpukake ing \olref{tab:prop-rules}.

\begin{table}
  \[\def\arraystretch{3}\begin{array}{|c|c|}
    \hline
    \AxiomC{\sFmla{\True}{\lnot !A}[\sigma]}
    \RightLabel{\TRule{\True}{\lnot}}
    \UnaryInfC{\sFmla{\False}{!A}[\sigma]}
    \DisplayProof
    &
    \AxiomC{\sFmla{\False}{\lnot !A}[\sigma]}
    \RightLabel{\TRule{\False}{\lnot}}
    \UnaryInfC{\sFmla{\True}{!A}[\sigma]}
    \DisplayProof
    \\[1ex]
    \hline
    \AxiomC{\sFmla{\True}{!A \land !B}[\sigma]}
    \RightLabel{\TRule{\True}{\land}}
    \UnaryInfC{\sFmla{\True}{!A}[\sigma]}
    \noLine
    \UnaryInfC{\sFmla{\True}{!B}[\sigma]}
    \DisplayProof
    &
    \AxiomC{\sFmla{\False}{!A \land !B}[\sigma]}
    \RightLabel{\TRule{\False}{\land}}
    \UnaryInfC{$\sFmla{\False}{!A}[\sigma] \quad \mid \quad
      \sFmla{\False}{!B}[\sigma]$}
    \DisplayProof
    \\[2ex]
    \hline
    \AxiomC{\sFmla{\True}{!A \lor !B}[\sigma]}
    \RightLabel{\TRule{\True}{\lor}}
    \UnaryInfC{$\sFmla{\True}{!A}[\sigma] \quad \mid \quad
      \sFmla{\True}{!B}[\sigma]$}
    \DisplayProof
    &
    \AxiomC{\sFmla{\False}{!A \lor !B}[\sigma]}
    \RightLabel{\TRule{\False}{\lor}}
    \UnaryInfC{\sFmla{\False}{!A}[\sigma]}
    \noLine
    \UnaryInfC{\sFmla{\False}{!B}[\sigma]}
    \DisplayProof
    \\[2ex]
    \hline
    \AxiomC{\sFmla{\True}{!A \lif !B}[\sigma]}
    \RightLabel{\TRule{\True}{\lif}}
    \UnaryInfC{$\sFmla{\False}{!A}[\sigma] \quad \mid
      \quad \sFmla{\True}{!B}[\sigma]$}
    \DisplayProof
    &
    \AxiomC{\sFmla{\False}{!A \lif !B}[\sigma]}
    \RightLabel{\TRule{\False}{\lif}}
    \UnaryInfC{\sFmla{\True}{!A}[\sigma]}
    \noLine
    \UnaryInfC{\sFmla{\False}{!B}[\sigma]}
    \DisplayProof
    \\[2ex]
    \hline
  \end{array}\]
  \caption{Aturan !!{tableau} mawa prefiks kanggo panggandheng
    proposisional}
  \ollabel{tab:prop-rules}
\end{table}

Sarat katutupan padha karo kang ana ing !!{tableau}s lumrah, nanging
kita mbutuhake supaya ora mung !!{formula}s, nanging uga prefiks-prefikse
padha. Dadi sawijining cabang katutup yen ngemot loro-lorone
\[
\sFmla{\True}{!A}[\sigma] \quad\text{lan}\quad \sFmla{\False}{!A}[\sigma]
\]
kanggo sawijining prefiks $\sigma$ lan !!{formula}~$!A$.

Aturan kanggo netepake asumsi uga padha karo kang ana ing
!!{tableau}s lumrah, kejaba yen kanggo asumsi kita tansah nggunakake
prefiks~$1$. (Prefiks endi wae sejatine bisa dienggo, anggere padha
kanggo kabeh asumsi.) Dadi, upamane, kita nyebut yen
\[
!B_1, \dots, !B_n \Proves !A
\]
yen lan mung yen ana tableau katutup kanggo asumsi
\[
\sFmla{\True}{!B_1}[1], \dots, \sFmla{\True}{!B_n}[1],
\sFmla{\False}{!A}[1].
\]

Kanggo operator modal\iftag{prvBox}{\iftag{prvDiamond}{~$\Box$
    lan}{~$\Box$}}{}\iftag{prvDiamond}{~$\Diamond$}{}, prefiks ing
kesimpulan aturan kang ditrapake marang !!a{formula} kanthi
prefiks~$\sigma$ yaiku $\sigma.n$. Nanging, $n$ endi kang diidini
gumantung marang apa tandhane~$\True$ utawa~$\False$.

\iftag{prvBox}{Aturan $\TRule{\True}{\Box}$ ngambakake cabang
  kang ngemot $\sFmla{\True}{\Box !A}[\sigma]$ kanthi
  $\sFmla{\True}{!A}[\sigma.n]$.\iftag{prvDiamond}{ Mangkono uga,
    a}{}}{A}\iftag{prvDiamond}{turan $\TRule{\False}{\Diamond}$
  ngambakake cabang kang ngemot $\sFmla{\False}{\Diamond !A}[\sigma]$ kanthi
  $\sFmla{\False}{!A}[\sigma.n]$.}{}
\iftag{notprvBox,notprvDiamond}{Aturan iki}{Aturan-aturan iki} mung
bisa ditrapake kanggo prefiks~$\sigma.n$ kang \emph{wis} ana ing
cabang papan aturan iku ditrapake. Prefiks mangkono diarani ``wis
dienggo'' (ing cabang kasebut).

\iftag{prvBox}{Aturan $\TRule{\False}{\Box}$ ngambakake cabang
  kang ngemot $\sFmla{\False}{\Box !A}[\sigma]$ kanthi
  $\sFmla{\False}{!A}[\sigma.n]$.\iftag{prvDiamond}{ Mangkono uga,
    a}{}}{A}\iftag{prvDiamond}{turan $\TRule{\True}{\Diamond}$ ngambakake
  cabang kang ngemot $\sFmla{\True}{\Diamond !A}[\sigma]$ kanthi
  $\sFmla{\True}{!A}[\sigma.n]$.}{}
\iftag{notprvBox,notprvDiamond}{Aturan iki}{Aturan-aturan iki},
nanging, mung bisa ditrapake kanggo prefiks~$\sigma.n$ kang
\emph{durung} ana ing cabang papan aturan iku ditrapake. Prefiks
mangkono diarani ``anyar'' (tumrap cabang kasebut).

Aturan-aturan kasebut disajekake ing \olref{tab:rules-K}.

\begin{table}
  \begin{center}
    \def\arraystretch{3}\def\fCenter{}
    \begin{tabular}{|c|c|}
    \hline
    \iftag{prvBox}{
      \AxiomC{\sFmla{\True}{\Box !A}[\sigma]}
      \RightLabel{\TRule{\True}{\Box}}
      \UnaryInfC{\sFmla{\True}{!A}[\sigma.n]}
      \DisplayProof
      &
      \AxiomC{\sFmla{\False}{\Box !A}[\sigma]}
      \RightLabel{\TRule{\False}{\Box}}
      \UnaryInfC{\sFmla{\False}{!A}[\sigma.n]}
      \DisplayProof\\
      $\sigma.n$ wis dienggo & $\sigma.n$ anyar
      \\[1ex]
      \hline}{}
    \iftag{prvDiamond}{
      \AxiomC{\sFmla{\True}{\Diamond !A}[\sigma]}
      \RightLabel{\TRule{\True}{\Diamond}}
      \UnaryInfC{\sFmla{\True}{!A}[\sigma.n]}
      \DisplayProof
      &
      \AxiomC{\sFmla{\False}{\Diamond !A}[\sigma]}
      \RightLabel{\TRule{\False}{\Diamond}}
      \UnaryInfC{\sFmla{\False}{!A}[\sigma.n]}
      \DisplayProof\\
      $\sigma.n$ anyar & $\sigma.n$ wis dienggo
      \\[1ex]
      \hline}{}
    \end{tabular}
  \end{center}
  \caption{Aturan modal kanggo \Ax{K}.}
  \ollabel{tab:rules-K}
\end{table}

Sarat yen prefiks kanggo
\iftag{prvBox}{\TRule{\True}{\Box}}{\TRule{\False}{\Diamond}} kudu wis
dienggo iku perlu; tanpa sarat iki, kita bakal nganggep ing ngisor iki
minangka !!{tableau} katutup:
\iftag{notprvBox,notprvDiamond}{%
  \iftag{prvBox}{%
  \begin{oltableau}
    [\pFmla{\True}{\Box \formula{A}}{1}, just = \TAss
      [\pFmla{\False}{\lnot\Box\lnot \formula{A}}{1}, just = \TAss
        [\pFmla{\True}{\formula{A}}{1.1}, just = {\TRule{\True}{\Box}[1]}
          [\pFmla{\True}{\Box\lnot\formula{A}}{1},
            just ={\TRule{\False}{\lnot}[2]}
            [\pFmla{\True}{\lnot\formula{A}}{1.1},
              just = {\TRule{\True}{\Box}[4]}
              [\pFmla{\False}{\formula{A}}{1.1},
                  just ={\TRule{\True}{\lnot}[5]}, close]
            ]
          ]
        ]
      ]
    ]
  \end{oltableau}
  }{
  \begin{oltableau}
    [\pFmla{\True}{\lnot\Diamond\lnot \formula{A}}{1}, just = \TAss
      [\pFmla{\False}{\Diamond\formula{A}}{1}, just = \TAss
        [\pFmla{\False}{\formula{A}}{1.1}, just = {\TRule{\False}{\Diamond}[2]}
          [\pFmla{\False}{\Diamond\lnot\formula{A}}{1},
            just ={\TRule{\True}{\lnot}[1]}
            [\pFmla{\False}{\lnot\formula{A}}{1.1},
              just = {\TRule{\False}{\Diamond}[4]}
              [\pFmla{\True}{\formula{A}}{1.1},
                  just ={\TRule{\True}{\lnot}[5]}, close]
            ]
          ]
        ]
      ]
    ]
  \end{oltableau}
}}{
  \begin{oltableau}
    [\pFmla{\True}{\Box \formula{A}}{1}, just = \TAss
      [\pFmla{\False}{\Diamond \formula{A}}{1}, just = \TAss
        [\pFmla{\True}{\formula{A}}{1.1}, just = {\TRule{\True}{\Box}[1]}
          [\pFmla{\False}{\formula{A}}{1.1},
            just = {\TRule{\False}{\Diamond}[2]}, close]
        ]
      ]
    ]
\end{oltableau}}

Nanging $\Box \formula{A} \Entails/ \Diamond \formula{A}$, mula sistem
bukti kita bakal ora sah. Mangkono uga, $\Diamond \formula{A} \Entails/
\Box \formula{A}$, nanging tanpa sarat yen prefiks kanggo
\iftag{prvBox}{\TRule{\False}{\Box}}{\TRule{\True}{\Diamond}} kudu
anyar, ing ngisor iki bakal dadi tableau katutup: \iftag{notprvBox,notprvDiamond}{%
  \iftag{prvBox}{%
  \begin{oltableau}
    [\pFmla{\True}{\lnot\Box\lnot \formula{A}}{1}, just = \TAss
      [\pFmla{\False}{\Box \formula{A}}{1}, just = \TAss
        [\pFmla{\False}{\formula{A}}{1.1}, just = {\TRule{\True}{\Box}[2]}
          [\pFmla{\False}{\Box\lnot\formula{A}}{1},
            just ={\TRule{\True}{\lnot}[1]}
            [\pFmla{\False}{\lnot\formula{A}}{1.1},
              just = {\TRule{\False}{\Box}[4]}
              [\pFmla{\True}{\formula{A}}{1.1},
                  just ={\TRule{\False}{\lnot}[5]}, close]
            ]
          ]
        ]
      ]
    ]
  \end{oltableau}
  }{
  \begin{oltableau}
    [\pFmla{\True}{\Diamond \formula{A}}{1}, just = \TAss
      [\pFmla{\False}{\lnot\Diamond\lnot\formula{A}}{1}, just = \TAss
        [\pFmla{\True}{\formula{A}}{1.1}, just = {\TRule{\True}{\Diamond}[1]}
          [\pFmla{\True}{\Diamond\lnot\formula{A}}{1},
            just ={\TRule{\False}{\lnot}[2]}
            [\pFmla{\True}{\lnot\formula{A}}{1.1},
              just = {\TRule{\True}{\Diamond}[4]}
              [\pFmla{\False}{\formula{A}}{1.1},
                  just ={\TRule{\True}{\lnot}[5]}, close]
            ]
          ]
        ]
      ]
    ]
  \end{oltableau}
}}{
  \begin{oltableau}
    [\pFmla{\True}{\Diamond \formula{A}}{1}, just = \TAss,name=one
      [\pFmla{\False}{\Box \formula{A}}{1}, just = \TAss, name=two
        [\pFmla{\True}{\formula{A}}{1.1}, just = {\TRule{\True}{\Diamond}[1]}
          [\pFmla{\False}{\formula{A}}{1.1},
            just = {\TRule{\False}{\Box}[2]}, close]
        ]
      ]
    ]
  \end{oltableau}}

\end{document}
